\chapter{Modelagem e Implementação do Data Warehouse}
\label{cap:dw}

Este capítulo registra as decisões técnicas na construção do DW para a execução orçamentária de Juiz de Fora. A ordem de exposição acompanha o fluxo dos dados: fontes públicas, extração, \textit{staging}, modelagem dimensional e camada analítica que alimenta os \textit{dashboards}. O foco recai sobre escolhas e seus motivos, não sobre exaustividade de código. Trechos de implementação entram quando ajudam a fixar um ponto concreto.

\section{Visão Geral da Arquitetura}

O sistema se organiza em quatro camadas, mostradas na Figura~\ref{fig:arquitetura-dw}: extração, \textit{staging}, \textit{warehouse} dimensional e consumo via \textit{dashboards}. Separar fisicamente essas camadas responde a uma exigência operacional: cada uma tem função definida e pode ser reexecutada sem arrastar as demais. Coleta, regra de negócio e apresentação ficam desacopladas. Quando o portal muda o \textit{layout} de uma planilha, o ajuste tende a ficar confinado à extração e ao \textit{staging}, sem reescrever consultas analíticas.

\begin{figure}[ht]
\begin{center}
\begin{tabular}{|l|}
\hline
\textbf{Fontes} \\
Portal Transparência PJF (XLS/XLSX) $\cdot$ STN (Ementário XLSX) \\
\hline
$\downarrow$ \\
\hline
\textbf{Extração e Orquestração} (Dagster + Python) \\
\texttt{despesa\_mensal\_consolidada}, \texttt{receita\_mensal\_comparativa}, \\
\texttt{receita\_mensal\_prevista}, \texttt{stn\_ementario\_natureza\_receita} \\
\hline
$\downarrow$ \\
\hline
\textbf{Staging} (DuckDB, \textit{schema} \texttt{stg}, dbt views) \\
Deduplicação, normalização leve, separação por layout \\
\hline
$\downarrow$ \\
\hline
\textbf{Warehouse Dimensional} (DuckDB, \textit{schema} \texttt{dwh}, dbt tables) \\
8 dimensões + 3 tabelas de fatos (modelagem Kimball) \\
\hline
$\downarrow$ \\
\hline
\textbf{Consumo} (Streamlit + Plotly) \\
7 páginas analíticas com filtros globais \\
\hline
\end{tabular}
\caption{Camadas do Data Warehouse construído para análise da execução orçamentária de Juiz de Fora.}
\label{fig:arquitetura-dw}
\end{center}
\end{figure}

A arquitetura segue de perto o padrão \textit{medallion} (\textit{bronze}/\textit{silver}/\textit{gold}), com adaptações locais. Os arquivos brutos ficam em \texttt{data/}, fora do banco. Tabelas brutas e \textit{views} de \textit{staging} compartilham o \textit{schema} \texttt{stg}: Dagster escreve as primeiras, dbt gera as segundas. O \textit{warehouse} dimensional ocupa o \textit{schema} \texttt{dwh}. A divisão por \textit{schema} deixa visível a fronteira entre dado bruto e dado modelado. Alterar uma transformação dbt não obriga a baixar de novo planilhas do portal.

\section{Stack Tecnológico}

A pilha foi montada para rodar em máquina local, com \textit{deploy} simples e ferramentas já difundidas em engenharia de dados. Para um TCC com uma pessoa na implementação, isso pesa tanto quanto desempenho bruto: quem for reproduzir o trabalho precisa conseguir subir o ambiente sem contratar infraestrutura. A Tabela~\ref{tab:stack} resume os componentes.

\begin{table}[htb]
\caption{Pilha técnica utilizada no Data Warehouse}
\label{tab:stack}
\centering
\small
\begin{tabular}{p{3.0cm} p{2.0cm} p{8.0cm}}
\hline
\textbf{Componente} & \textbf{Versão} & \textbf{Função no projeto} \\
\hline
Python & 3.13+ & Linguagem de extração e código geral \\
Dagster & 1.12 & Orquestrador de \textit{assets} \\
\texttt{pandas} & 3.0 & Leitura e normalização de planilhas \\
DuckDB & 1.5 & Banco analítico embarcado, formato OLAP colunar \\
dbt-core & 1.11 & Modelagem dimensional em SQL \\
dbt-duckdb & 1.10 & Adaptador dbt para DuckDB \\
\texttt{dbt\_utils} & 1.3 & Macros para chaves \textit{surrogate} e \textit{date spine} \\
Streamlit & 1.57 & Construção dos \textit{dashboards} \\
Plotly & 6.7 & Visualizações interativas \\
pytest & 9.0 & Testes da camada de \textit{dashboard} \\
\textit{ruff} + \textit{sqlfluff} & n/a & \textit{Linting} de Python e SQL \\
\hline
\end{tabular}
\end{table}

DuckDB merece explicação porque desloca o projeto de soluções mais habituais em DW (PostgreSQL, BigQuery, Snowflake). Trata-se de banco analítico em processo, sem servidor, com armazenamento colunar voltado a agregações. Para escopo municipal e volume na casa de milhões de linhas, três fatores inclinaram a balança. Não há servidor para instalar: o banco é o arquivo \texttt{data/execucao\_orcamentaria.duckdb}. Consultas analíticas competem, neste porte, com opções pagas. A ponte com \texttt{pandas} e Arrow encurta o caminho entre Python e SQL. O preço é a ausência de concorrência multiusuário em produção, irrelevante enquanto o consumo for local ou de poucos analistas, mas decisivo se o sistema migrar para portal público com tráfego simultâneo.

Dagster e dbt dividem o trabalho de forma comum na engenharia de dados atual: o primeiro orquestra \textit{assets} de forma declarativa, o segundo concentra transformações em SQL. A integração via \texttt{dagster-dbt} amarra as duas partes. O \textit{asset} \texttt{dbt\_build} só roda depois dos \textit{assets} de extração, de modo que modelagem e carga chegam juntas a cada ciclo.

\section{Extração: Portal da Transparência e STN}
\label{sec:extracao}

O Portal da Transparência da Prefeitura de Juiz de Fora não expõe API REST. A coleta depende de \texttt{HTTP GET} sobre planilhas Excel em URLs previsíveis. Esse desenho é típico de portais municipais brasileiros: a transparência existe, mas no formato que a prefeitura já usa internamente, não no formato que um pipeline automatizado preferiria. Daí o padrão adotado: baixar arquivo, localizar cabeçalho, interpretar com \texttt{pandas}. A Tabela~\ref{tab:fontes} sintetiza as fontes consumidas.

\begin{table}[htb]
\caption{Fontes de dados consumidas pelo \textit{pipeline} de extração}
\label{tab:fontes}
\centering
\small
\begin{tabular*}{\textwidth}{@{\extracolsep{\fill}}p{2.3cm}cc p{8.8cm}@{}}
\hline
\textbf{\textit{Asset}} & \textbf{Origem} & \textbf{Partição} & \textbf{Conteúdo} \\
\hline
\begin{minipage}[t]{2.1cm}\raggedright\path{despesa_mensal_consolidada}\end{minipage} & PJF & Mensal & Empenhos, liquidações e pagamentos do mês \\
\begin{minipage}[t]{2.1cm}\raggedright\path{receita_mensal_comparativa}\end{minipage} & PJF & Mensal & Previsão e arrecadação acumuladas no mês \\
\begin{minipage}[t]{2.1cm}\raggedright\path{receita_mensal_prevista}\end{minipage} & PJF & Anual & Previsão mensal de receitas para o exercício \\
\begin{minipage}[t]{2.1cm}\raggedright\path{stn_ementario_natureza_receita}\end{minipage} & STN & Anual & Códigos e descrições oficiais de natureza de receita \\
\hline
\end{tabular*}
\end{table}

Três decisões de extração ilustram como o pipeline reage a problemas reais de fonte pública, não a um \textit{dataset} estável de laboratório.

\subsection{Detecção Dinâmica de Cabeçalho}

As planilhas de despesa não fixam a linha de cabeçalho. O arquivo abre com metadados (logotipo, título, período) antes da tabela. A extração percorre as primeiras linhas até achar a célula \texttt{"Unidade Administrativa"} e só então define o início dos dados. Pequenas mudanças de \textit{layout} entre meses deixam de quebrar a carga.

\subsection{Quebra de \textit{Layout} em 2025-05}

A planilha de receita comparativa mudou de formato no arquivo de maio de 2025 (partição \texttt{2505}). Em vez de descartar histórico ou forçar um parser único, o código mantém dois leitores: \texttt{read\_receita\_comparativa\_pre2505} e \texttt{read\_receita\_comparativa\_2505}. A escolha ocorre na extração, pela chave de partição. Anos anteriores continuam reprocessáveis sem perda.

\subsection{Idempotência por Estratégia de Escrita}

Cada \textit{asset} usa a estratégia de escrita que combina com sua semântica. Despesa e receita comparativa entram por \textit{append}; a deduplicação fica no \textit{staging}, por \texttt{nm\_arquivo} e \texttt{dt\_atualizacao}. O ementário da STN, anual e canônico, usa \textit{overwrite by partition}, substituindo todas as linhas do ano a cada execução. Os caminhos diferem, mas o efeito desejado é o mesmo: reexecutar o pipeline para a mesma partição devolve o mesmo estado final do \textit{warehouse}.

\section{Camada de Staging}

O \textit{staging} aproxima o bruto do formato que as regras de negócio exigem, sem aplicar regra de negócio ainda. Há deduplicação, remoção de linhas agregadas (\texttt{TOTAIS GERAIS}), normalização de códigos de natureza (só dígitos, sem sufixos como \texttt{.0}) e \textit{trim} de espaços. Quatro \textit{views} dbt cobrem essa etapa, uma por fonte.

A deduplicação de despesa mostra por que \textit{append} na extração exige cuidado a jusante. O mesmo arquivo pode entrar várias vezes, sobretudo em \textit{backfills}. A \textit{view} de \textit{staging} resolve com janela:

\begin{singlespacing}
\begin{lstlisting}[language=SQL,basicstyle=\ttfamily\small,frame=single,breaklines=true]
select *
from {{ source('stg', 'pjf_despesa_mensal_consolidada') }}
qualify rank() over (
    partition by nm_arquivo
    order by dt_atualizacao desc
) = 1
\end{lstlisting}
\end{singlespacing}

A cláusula \texttt{qualify}, suportada pelo DuckDB e por dialetos SQL recentes, filtra o resultado de funções de janela sem subconsulta. Para cada arquivo de origem, permanece só a versão mais recente carregada.

\section{Modelo Dimensional}

A modelagem segue os quatro passos de Kimball: processos de negócio, granularidade, dimensões e fatos. Dois processos foram tratados em separado: \textbf{execução de despesa} e \textbf{execução de receita}. Unificá-los num único fato misturaria grãos, dimensões e calendários de publicação distintos.

\subsection{Tabelas de Fatos}

O modelo conta com três tabelas de fatos, descritas na Tabela~\ref{tab:fatos}.

\begin{table}[htb]
\caption{Tabelas de fatos do \textit{warehouse}}
\label{tab:fatos}
\centering
\small
\begin{tabular*}{\textwidth}{@{\extracolsep{\fill}}p{2.2cm}p{7.2cm}p{5.2cm}@{}}
\hline
\textbf{Tabela} & \textbf{Granularidade} & \textbf{Medidas} \\
\hline
\begin{minipage}[t]{2.0cm}\raggedright\path{fct_despesa}\end{minipage} & Linha do arquivo mensal consolidado (combinação de empenho, unidade, fornecedor, funcional, natureza, fonte e datas) & \path{vl_empenhado_mes}, \path{vl_liquidado_mes}, \path{vl_pago_mes} \\
\begin{minipage}[t]{2.0cm}\raggedright\path{fct_receita}\end{minipage} & Natureza de receita por mês de referência & \path{vl_previsto_mensal}, \path{vl_arrecadada_mes} \\
\begin{minipage}[t]{2.0cm}\raggedright\path{fct_receita_acumulada}\end{minipage} & Natureza de receita por mês de referência & \path{vl_previsao_inicial_comparativa}, \path{vl_previsao_atualizada}, \path{vl_arrecadada_ano}, \path{vl_a_realizar} \\
\hline
\end{tabular*}
\end{table}

A granularidade de \texttt{fct\_despesa} exige cuidado. No ideal kimballiano, cada linha seria um empenho ou um estágio (empenho, liquidação, pagamento). O portal, porém, publica relatórios mensais já consolidados: cada linha agrega valores por combinação de dimensões. Modelar nesse grão respeita o que a fonte entrega, sem simular atomicidade inexistente. A troca é clara: perguntas sobre transações individuais (um empenho, um pagamento) ficam fora do alcance. Isso vem do portal, não de falha do modelo.

\texttt{fct\_receita} e \texttt{fct\_receita\_acumulada} também nascem de semânticas diferentes. O arquivo de previsão mensal traz valor previsto por mês, alimentando \texttt{fct\_receita} após \textit{unpivot}. O comparativo traz acumulados no ano (previsão inicial, previsão atualizada, arrecadada, a realizar), alimentando \texttt{fct\_receita\_acumulada}. Juntar tudo numa tabela única exigiria misturar medidas de ritmo mensal com medidas de saldo anual.

\subsection{Dimensões}

O modelo inclui oito dimensões, descritas na Tabela~\ref{tab:dimensoes}.

\begin{table}[htb]
\caption{Dimensões do \textit{warehouse}}
\label{tab:dimensoes}
\centering
\small
\begin{tabular}{p{4.5cm} p{8.5cm}}
\hline
\textbf{Dimensão} & \textbf{Conteúdo principal} \\
\hline
\texttt{dim\_tempo} & Dia, mês, trimestre, semestre, ano (\textit{date spine} de 2014 até 1 ano à frente) \\
\texttt{dim\_unidade\_administrativa} & Órgão executor da despesa \\
\texttt{dim\_funcional\_pragmatica} & Código e descrição da funcional programática \\
\texttt{dim\_natureza\_despeza} & Código e descrição da natureza de despesa \\
\texttt{dim\_fonte\_recurso} & Código e descrição da fonte de recurso \\
\texttt{dim\_fornecedor} & CPF/CNPJ e nome do fornecedor \\
\texttt{dim\_empenho} & Número da nota, modalidade, licitação, referência legal, processo e descrição \\
\texttt{dim\_natureza\_receita} & Código e descrição da natureza de receita (enriquecida pelo ementário STN) \\
\hline
\end{tabular}
\end{table}

\texttt{dim\_natureza\_receita} cruza duas origens. Quando o ementário da STN cobre o código, prevalece a descrição oficial; nos demais casos, entra a descrição publicada pela PJF. Quem consulta o \textit{dashboard} vê nomenclatura nacional padronizada onde ela existe, sem apagar códigos que só aparecem no município.

O nome \texttt{dim\_natureza\_despeza} mantém um erro de digitação do projeto. Registro aqui por transparência e como aviso: renomear modelos dbt exige revisar referências, testes e consultas do \textit{dashboard}.

\subsection{Chaves \textit{Surrogate}}

Todas as dimensões usam chaves \textit{surrogate} geradas por \texttt{dbt\_utils.generate\_surrogate\_key}, com MD5 sobre concatenação determinística dos campos naturais. Kimball recomenda esse padrão por dois motivos operacionais. A chave não muda se um rótulo de exibição for corrigido na fonte. A junção entre fatos volumosos e dimensões usa um único \texttt{varchar(32)}. O hash determinístico garante que reprocessar o pipeline recria as mesmas chaves, preservando idempotência.

\subsection{Esquema Estrela}

As ligações entre fatos e dimensões formam três estrelas que compartilham \texttt{dim\_tempo}. A Figura~\ref{fig:estrelas} representa essas conexões.

\begin{figure}[ht]
\begin{center}
\begin{tabular}{|l|}
\hline
\textbf{Estrela de Despesa} \\
\texttt{fct\_despesa} $\leftarrow$ dim\_tempo (empenho, liquidação) \\
\texttt{fct\_despesa} $\leftarrow$ dim\_unidade\_administrativa \\
\texttt{fct\_despesa} $\leftarrow$ dim\_fornecedor \\
\texttt{fct\_despesa} $\leftarrow$ dim\_funcional\_pragmatica \\
\texttt{fct\_despesa} $\leftarrow$ dim\_natureza\_despeza \\
\texttt{fct\_despesa} $\leftarrow$ dim\_fonte\_recurso \\
\texttt{fct\_despesa} $\leftarrow$ dim\_empenho \\
\hline
\textbf{Estrela de Receita} \\
\texttt{fct\_receita} $\leftarrow$ dim\_tempo, dim\_natureza\_receita \\
\hline
\textbf{Estrela de Receita Acumulada} \\
\texttt{fct\_receita\_acumulada} $\leftarrow$ dim\_tempo, dim\_natureza\_receita \\
\hline
\end{tabular}
\caption{Esquemas estrela do \textit{warehouse}, compartilhando \texttt{dim\_tempo}.}
\label{fig:estrelas}
\end{center}
\end{figure}

Nenhuma dimensão implementa SCD (\textit{slowly changing dimension}). Todas são tipo 1: guardam só o estado mais recente da fonte. Para descrições de natureza no padrão brasileiro, isso costuma bastar, porque a STN revisa ementários em ciclos anuais e as mudanças são pontuais. Reorganizações administrativas do município, mais raras e mais abruptas, ficariam mal registradas. Capturá-las exigiria SCD tipo 2, tarefa deixada para evolução futura.

\section{Processos de ETL e Idempotência}

A unidade de idempotência é a partição temporal: mês de competência na PJF, ano de referência na STN. Reexecutar o pipeline para uma partição deve devolver o mesmo \textit{warehouse}, independentemente de quantas cargas anteriores existam. Três mecanismos sustentam essa propriedade.

A estratégia de escrita por fonte (Seção~\ref{sec:extracao}) trata STN com sobrescrita por partição e PJF com \textit{append} mais deduplicação.

As chaves \textit{surrogate} determinísticas evitam linhas espúrias em dimensões quando uma partição é recarregada: o mesmo conjunto de campos naturais gera o mesmo MD5.

A materialização \textit{table}, e não incremental, nas tabelas de fatos faz o dbt reconstruir o conjunto inteiro a partir do \textit{staging} a cada execução. Para o volume municipal atual, o \textit{full refresh} elimina erros típicos de \textit{merge} incremental mal definido. Se o dado crescer, modelos incrementais com chave única estável podem substituir essa escolha sem mudar o contrato dos \textit{dashboards}.

\section{Qualidade de Dados e Testes}

Os testes operam em três níveis. No dbt, arquivos \texttt{\_schema.yml} checam unicidade e nulidade de chaves de dimensão, unicidade das combinações de \textit{surrogate keys} nos fatos, presença de medidas e integridade referencial. Em regras de negócio, testes singulares em SQL verificam invariantes de receita: arrecadação acumulada não pode ser menor que o acumulado dos meses anteriores (monotonicidade no ano), e o valor a realizar deve bater com previsão atualizada menos arrecadado no ano. No \textit{dashboard}, \texttt{pytest} cobre formatadores, classificação de natureza e função de governo, além do \textit{schema} esperado das tabelas do \textit{warehouse}.

Em dados fiscais publicados, falha silenciosa custa caro: um valor negativo por erro de \textit{parsing} ou uma quebra de monotonicidade entre meses pode virar conclusão errada sobre execução orçamentária. O conjunto de testes não cobre tudo, mas ataca as falhas mais prováveis: chave duplicada em dimensão, medida monetária inválida, carga parcial entre meses e \textit{drift} de \textit{schema} após mudança de \textit{layout} na fonte. A Tabela~\ref{tab:dq} resume os testes singulares.

\begin{table}[htb]
\caption{Testes singulares de qualidade de dados}
\label{tab:dq}
\centering
\small
\begin{tabular}{p{6.5cm} p{6.5cm}}
\hline
\textbf{Teste} & \textbf{Invariante verificada} \\
\hline
\texttt{dq\_fct\_receita\_acumulada\_nao\_menor\_que\_mes} & Valor acumulado no ano deve ser maior ou igual à arrecadação no mês corrente \\
\texttt{dq\_fct\_receita\_acumulada\_monotonicidade} & Valor acumulado no ano deve ser não decrescente mês a mês, por natureza \\
\texttt{dq\_fct\_receita\_vl\_a\_realizar} & Valor a realizar deve coincidir com previsão atualizada menos arrecadado no ano \\
\hline
\end{tabular}
\end{table}

\section{Limitações da Implementação}

Registrar o que o sistema não faz é tão importante quanto descrever o que faz. Três limitações marcam o recorte atual e orientam quem quiser reproduzir ou estender o trabalho.

\subsection{Granularidade da Despesa}

A fonte publica relatórios mensais consolidados, não transações individuais. Análises empenho a empenho exigiriam outras bases (sistemas internos da prefeitura, eventualmente obtidos por LAI). O modelo define o grão pela linha do arquivo, mas quem lê \texttt{fct\_despesa} precisa saber que os valores são movimentação mensal por combinação de dimensões, não atos administrativos isolados.

\subsection{Ausência de Orçamento Autorizado}

O \textit{warehouse} traz empenhado, liquidado e pago, mas não a dotação autorizada na LOA. Sem autorizado, indicadores como percentual de execução sobre o previsto na lei ficam indisponíveis. Incluir a LOA implica nova coleta (muitas vezes PDF) e parsing próprio, previstos como continuidade.

\subsection{Classificação MCASP Funcional Incompleta}

A dimensão \texttt{dim\_funcional\_pragmatica} guarda o código composto da funcional programática como \textit{string} única, sem separar função e subfunção. O \textit{dashboard} compensa com dicionário local de funções de governo, o que basta para cortes por área (saúde, educação, segurança), mas fica aquém da hierarquia MCASP completa. Uma \texttt{dim\_funcao} alimentada pelo ementário oficial seria o passo seguinte natural.

Essas limitações delimitam perguntas respondíveis hoje e apontam extensões concretas: nova fonte para dotação, grão mais fino de despesa, dimensão funcional desagregada. O DW entregue cobre execução publicada no portal; não substitui o ERP municipal nem fecha sozinho o ciclo orçamentário completo.
